\documentclass[11pt]{article}
\usepackage{icmlko}
\begin{document}

\title{전체를 껐을 때의 지도는 부호마저 반대였다 --- 두 도메인이 서로 뒤집히는 것을 보고서야 한 도메인을 고칠 수 있었다}
\author{월드모델 실험실}
\date{노트 563 · 564 · 2026년 8월 5일}
\maketitle

\begin{abstract}
노트 $559$ 가 검색 축의 \texttt{zero\_is\_data} 를 전역으로 바꾸려다
기각됐고, 그때 낸 \textbf{도메인별 표}에서 도서가 $+0.0459$ 로 제일 컸다.
노트 $560$ 이 도서만 바꿔 보니 $\mathbf{-0.0820}$ --- 부호가 뒤집혔다.
그래서 반대쪽도 봤다.

\begin{center}
\begin{tabular}{lrr}
\hline
 & 전역 변경의 도메인별 차 & \textbf{그 도메인만 바꿈} \\
\hline
도서 & $+0.0459$ & $\mathbf{-0.0820}$ \\
\textbf{게임} & $-0.0137$ & $\mathbf{+0.0138}$ \\
\hline
\end{tabular}
\end{center}

\textbf{둘 다 뒤집혔고, 서로 반대 방향이다.} 전역 제거의 도메인별 표는
크기가 부정확한 것이 아니라 \textbf{부호를 못 쓴다.}
\end{abstract}

\begin{figure}[t]
\centering
\includegraphics[width=\textwidth]{figs/f1.pdf}
\caption{\textbf{가} --- 부호가 뒤집힌다. \textbf{나} --- 관측-$0$ 의 정체.
\textbf{다} --- 채택.}
\end{figure}

\section{게임의 관측-$0$ 은 수요가 아니다}

\begin{center}
\begin{tabular}{lrrr}
\hline
게임 & $n$ & 라벨 중앙 & \textbf{한글 제목} \\
\hline
값 있음 & $305$ & $8.64$ & $36\%$ \\
\textbf{관측-$0$} & $42$ & $7.74$ & $\mathbf{100\%}$ \\
안 긁힘 & $135$ & $8.12$ & $6\%$ \\
\hline
\end{tabular}
\end{center}

\textbf{영문 제목 게임 중 관측-$0$ 은 $0$ 건이다.} 영문(글로벌 대작)은
긁으면 늘 값이 나오고, 관측-$0$ 은 \emph{전부} 한글 제목이다.

\begin{quote}
그러면 그 마스크는 수요가 아니라 \textbf{``한글 제목 소규모 국산인가''}라는
도메인 안 하위집단 표지다. 값 $0$(백분위 최저)으로 두면 ``검색량이 제일
낮다''는 \textbf{잘못된 크기}를 주고, 결측으로 두면 ``모른다''가 되어 낫다.
\end{quote}

\textbf{도서는 반대다} --- 값있음 $97$ 행과 관측-$0$ $109$ 행이 \emph{둘 다
한글 $100\%$} 다. 언어로 안 갈리므로 관측-$0$ 이 진짜 ``무명하다''이고,
그래서 빼면 손해다($-0.0820$).

\section{채택 --- 게임만}

\begin{center}
\begin{tabular}{lrrl}
\hline
자 (짝 씨앗 $40$) & 전 & 후 & 차 \\
\hline
판 & $+0.4673$ & $+0.4678$ & $+0.0005$ ($23/40$, 문다) \\
\textbf{날짜 통제 판} & $+0.5047$ & $+0.5054$ & $\mathbf{+0.0007}$ ($0$ 밖 양수) \\
무리 안 판 & $+0.3860$ & $+0.3869$ & $+0.0009$ (문다) \\
KR 만화 & $+0.6378$ & $+0.6380$ & $+0.0001$ (문다) \\
비게임 앱 & $+0.2901$ & $+0.2895$ & $-0.0006$ (문다) \\
\hline
\textbf{게임} & $+0.6459$ & $+0.6597$ & $\mathbf{+0.0138}$ ($\mathbf{40/40}$) \\
\hline
\end{tabular}
\end{center}

대조(씨앗 $12$): 게임 검색을 \textbf{통째로} 가리면 $-0.0587$ --- 값은
필요하고 \textbf{관측-$0$ 만} 해롭다. \textbf{판은 안 움직인다} --- 노트
$351 \cdot 553$ 과 같은 모양이다.

\section{안쪽 검증이 전역 표를 이겼다}

노트 $562$ 가 도메인별 모형 고르기를 기각했다 --- 배정이 무작위와 구별이
안 됐다($C-D = -0.0003$). \textbf{그런데 게임만 씨앗 $12$ 중 $12$ 번 Q 를
골랐다}(안쪽 검증 $100$ 행). 만장일치는 셋뿐이었고 여덟은 흔들렸다.

\begin{quote}
\textbf{그 만장일치가 맞았다.} 전역 표가 게임을 $-0.0137$ 로 적을 때 안쪽
검증은 $+$ 를 가리켰고, 단일 전환은 $+0.0138$ 이었다. \textbf{고르기가
전체적으로 안 되는 것과 한 도메인에서 맞는 것은 다른 이야기}이고, 그 경계는
노트 $562$ 의 크기 법칙($n > (2/e)^2$)이 정한다.
\end{quote}

\section*{내가 틀린 것 --- 진단에 다른 라벨을 썼다}

기제를 \texttt{harness.load()} 라벨로 쟀는데, \textbf{챔피언은
\texttt{LABEL\_FIX} 로 게임 라벨을 \texttt{y\_w30}(출시 $30$ 일 리뷰)로
바꾼다}(노트 $210$).

\begin{center}
\begin{tabular}{lrr}
\hline
 & \texttt{load()} & \textbf{챔피언} \\
\hline
``관측-$0$ 인가''$\leftrightarrow$라벨 & $-0.1721$ & $\mathbf{-0.1157}$ \\
``한글 제목''$\leftrightarrow$라벨 & $-0.2009$ & $\mathbf{-0.0840}$ \\
관측-$0$ 대 안긁힘 & $p=10^{-4}$ & $p=0.53$ \\
\hline
\end{tabular}
\end{center}

\textbf{한글 제목 구성은 안 바뀐다}(라벨과 무관) --- 값있음 $36\%$ ·
관측-$0$ $\mathbf{100\%}$ · 안긁힘 $6\%$. 기제의 핵심은 그것이므로 결론은
선다.

\begin{quote}
\textbf{조항: 도메인 진단은 챔피언 자료로 한다.} \texttt{harness.load()} 는
\texttt{fixaxes.apply} $\cdot$ \texttt{apply\_labels} 를 안 거치므로 게임
라벨이 다르고 차단 축도 살아 있다. 노트 $435$ 가 ``원배열이 아니라 모형이
보는 설계행렬에서 확인한다''고 적은 것과 같은 자리이고, \textbf{이번엔 라벨
쪽}이었다.
\end{quote}

\section*{조항}

\begin{quote}
\textbf{전역 변경의 도메인별 표는 부호도 못 쓴다.} 노트 $551$ 이 ``크기가
섞인다''로 세운 조항을 두 사례가 \textbf{``부호가 뒤집힌다''}로 올린다
(도서 $+ \to -$ · 게임 $- \to +$). 그 표는 \emph{어느 도메인을 볼지}
고르는 데만 쓰고, \textbf{부호와 크기는 단일 전환으로 다시 잰다.}
\end{quote}

\section*{남은 것}

\begin{itemize}
\item \textbf{모바일의 ``제목이 한글인가''} --- \texttt{tagaxes} 채택 검사
셋을 다 통과한다(① 학습 $+0.4362$ · ② 다섯 조각 부호 전부 $+$ · ③ 최대
겹침 $0.371 < 0.39$). 모바일 무게 $13.1\%$. \textbf{이번 사이클에서 처음
나온 새 축 후보}다.
\item \textbf{게임에서는 같은 축이 검사 ①②를 못 넘는다}(학습 $-0.0489$ ·
부호 $[+,+,-,-,-]$). 두 도메인에서 부호가 반대인 것이 아니라 \emph{게임
쪽이 약하다}.
\item \textbf{도서 $\cdot$ 웹툰 $\cdot$ 모바일의 관측-$0$} --- 게임과 같은
검사를 아직 안 걸었다(``관측-$0$ 이 어떤 하위집단인가'').
\end{itemize}

\end{document}
